\documentclass{article}
\usepackage{graphicx}
\usepackage{booktabs}
\usepackage{vouch}

\title{When Does a Neighbour Beat a Line?\\[2pt] \large A toy learning-curve study}
\author{The vouch tutorial}
\date{}

\begin{document}
\maketitle

\begin{abstract}
We compare a linear model with $k$-nearest neighbours as the training set grows.
The linear model is better with little data, but $k$-NN has the higher accuracy
from \vouch{crossover} training points on and finishes \vouch{gap} points ahead.
\end{abstract}

\section{Setup}

Points are drawn uniformly from the square $[-1, 1]^2$ and labelled by which side
of a sine curve they fall on. We flip \vouch[.0pct]{experiment.param.noise} of the
labels at random, so no classifier can exceed \vouch{best_possible} test accuracy.
We train logistic regression and a \vouch{experiment.param.k}-nearest-neighbour
classifier on \vouch{experiment.param.sizes} points, test each on
\vouch{experiment.param.n_test} fresh points, and repeat every configuration with
\vouch{experiment.param.seeds} seeds.

\section{Results}

Figure~\ref{fig:curve} and Table~\ref{tab:curve} show test accuracy as the
training set grows. \vouchclaim{linear_wins_small}{With the fewest training points
the linear model is better}: \vouch{evaluate.linear.n_train_20.acc} against
\vouch{evaluate.knn.n_train_20.acc} for $k$-NN, whose vote is then dominated by
distant points. \vouchclaim{knn_wins_large}{With the most, $k$-NN is better}: it
reaches \vouch{evaluate.knn.n_train_640.acc}, against
\vouch{evaluate.linear.n_train_640.acc} for the linear model, and varies by only
\vouch{evaluate.knn.n_train_640.acc.std} across seeds.

The linear model lacks flexibility, not data: with the most training points its
training accuracy, \vouch{evaluate.linear.n_train_640.train_acc}, is no higher than
its test accuracy. A straight line cannot follow the sine boundary, while $k$-NN
can, and keeps improving towards the \vouch{best_possible} ceiling.

\begin{figure}[t]
  \centering
  \includegraphics[width=0.75\linewidth]{figures/learning_curve.pdf}
  \caption{Test accuracy against training set size, mean $\pm$ std over
    \vouch{experiment.param.seeds} seeds. The dashed line is the best possible
    accuracy, \vouch{best_possible}.}
  \label{fig:curve}
\end{figure}

\begin{table}[t]
  \centering
  \caption{Test accuracy by training set size, mean $\pm$ std over
    \vouch{experiment.param.seeds} seeds. The gap is $k$-NN minus the linear
    model, in percentage points.}
  \label{tab:curve}
  \begin{tabular}{rccr}
    \toprule
    Training points & Linear & $k$-NN & Gap \\
    \midrule
    \vouchtable{learning_curve}
    \bottomrule
  \end{tabular}
\end{table}

\end{document}
